\section{Second out-neighborhoods and a research problem}\label{ch07:sec:seymour}
This section states a problem about directed graphs that resisted every attempt at a proof for more than thirty years. Stating it exactly needs one more idea: distance measured along the arrows.

The deletion argument of Lemma~\ref{ch07:lem:shortest-walk} works just as well for directed walks: cutting out the part between two visits to the same vertex leaves a shorter directed walk with the same start and end, because the arrow that used to leave the second visit now leaves the first, which is the same vertex. So, as in the lemma, whenever there is a directed walk from $v$ to $w$, there is a shortest one, and it is a directed path. Its length is called the \emph{directed distance} from $v$ to $w$. Unlike distance in a graph, directed distance depends on the order of the two vertices. In the chain $a\to b\to c$, the directed distance from $a$ to $c$ is $2$, but there is no directed walk from $c$ back to $a$ at all.

Now let $v$ be a vertex of a directed graph with vertex set $V$. Write $N_1^+(v)$ for the set of out-neighbors of $v$, and call it the \emph{first out-neighborhood} of $v$. Its members are exactly the vertices at directed distance $1$ from $v$. The \emph{second out-neighborhood} of $v$ is the set
\[
N_2^+(v)=\{w\in V:\ w\ne v,\ w\notin N_1^+(v),\ \text{and }v\to u\to w\text{ for some vertex }u\}.
\]
The superscript $+$ is a reminder that we follow the arrows outward from $v$, and the subscript counts the steps. In words, $w$ belongs to $N_2^+(v)$ when some two-step walk $v\to u\to w$ reaches it but no shorter walk does: $w$ is not $v$ itself, at distance $0$, and not a first out-neighbor, at distance $1$. So $N_2^+(v)$ is the set of vertices at directed distance exactly two from $v$. It is a set of vertices, not of routes: if several intermediate vertices $u$ lead to the same $w$, then $w$ is still a single member of $N_2^+(v)$ and is counted once in $|N_2^+(v)|$.

In an oriented graph, the condition $w\ne v$ holds automatically, because a route $v\to u\to v$ would need edges in both directions between $v$ and $u$. In a directed graph that is not oriented, such a route can occur, and the condition $w\ne v$ is what prevents $v$ from counting as its own second out-neighbor.

Two small examples show the definition at work. In the oriented graph with edges $a\to b$, $b\to c$, $a\to c$, and $c\to d$, the two-step routes from $a$ are $a\to b\to c$ and $a\to c\to d$. The vertex $c$ is at the end of the first route, but it is not in $N_2^+(a)$, because the edge $a\to c$ already makes it a first out-neighbor. So
\[
N_1^+(a)=\{b,c\},\qquad N_2^+(a)=\{d\},\qquad |N_1^+(a)|=2,\qquad |N_2^+(a)|=1.
\]
In the directed cycle $a\to b\to c\to a$, each vertex has exactly one first and one second out-neighbor; for instance, $N_1^+(a)=\{b\}$ and $N_2^+(a)=\{c\}$. Here nothing is excluded: the only two-step route from $a$ is $a\to b\to c$, and $c$ is neither $a$ nor an out-neighbor of $a$.

In 1990 Paul Seymour proposed the following statement, now known as \emph{Seymour's second-neighborhood conjecture}:
\begin{quote}
Every nonempty finite oriented graph has a vertex $v$ with $|N_2^+(v)|\ge|N_1^+(v)|$.
\end{quote}
In words: some vertex reaches at least as many new vertices in its second step as in its first. In the four-vertex example above, the vertex $a$ fails the inequality, since $1<2$. The vertex $b$ satisfies it, since $N_1^+(b)=\{c\}$ and $N_2^+(b)=\{d\}$, so $1\ge1$. So does $d$, for a simpler reason. It is a sink, and a sink always satisfies the inequality: it has no out-neighbors, so it has no two-step routes either, and both of its out-neighborhoods are empty. So the interesting graphs are those in which every vertex has at least one out-neighbor.

The word ``oriented'' cannot be dropped from the conjecture. Consider the directed graph with two vertices $u$ and $v$ and both edges $u\to v$ and $v\to u$. The only two-step route from $u$ is $u\to v\to u$, which returns to $u$ itself, so $N_1^+(u)=\{v\}$ and $N_2^+(u)=\varnothing$. The same holds for $v$, so neither vertex satisfies the inequality.

For decades the conjecture remained open, although special cases were settled; the case of tournaments, for example, was proved by David Fisher in 1996. In September 2026, a \emph{preprint}, that is, a research paper made public before a journal has reviewed it, claimed a proof of the full conjecture. It was released in OpenAI's public mathematics repository together with a formal proof in the Lean proof assistant~\cite{oai173}. The repository's description of the formalized statement uses the same conventions as this section: the second out-neighborhood consists of the vertices at directed distance exactly two, so $v$ and its first out-neighbors are excluded, and graphs with sinks are allowed. The proof is long and draws on ideas well beyond this course. Section~\ref{ch16:sec:second} shows how the definition of $N_2^+(v)$ is written in Lean, and Section~\ref{ch17:sec:research} describes how to begin reading a formalized proof of this kind.

If a news report or an AI assistant tells you that ``Seymour's conjecture has been proved,'' the first thing to check is \emph{which statement} was proved. A proof that counted two-step routes instead of vertices, or that allowed first out-neighbors to count as second out-neighbors, would establish a weaker statement. Each change can only make the second count larger, and a larger second count makes the inequality easier to satisfy. The definitions in this section are what let you ask that question precisely.

We can already prove one special case completely.

\par\noindent\begin{minipage}{\linewidth}
\begin{proposition}\label{ch07:prop:acyclic-seymour}
Every nonempty finite acyclic oriented graph has a vertex $v$ with $|N_2^+(v)|\ge|N_1^+(v)|$.
\end{proposition}
\end{minipage}\par

\begin{proof}
By Theorem~\ref{thm:sink}, the graph has a sink $v$. As we saw above, both out-neighborhoods of a sink are empty. Hence $|N_2^+(v)|=0\ge0=|N_1^+(v)|$.
\end{proof}

This is a complete proof for acyclic graphs, and no proof at all for graphs with directed cycles, which is where all the difficulty lies. In fact, a counterexample to the conjecture would have to be a nonempty finite oriented graph without a sink, and by Theorem~\ref{thm:sink} such a graph must contain a directed cycle. Keeping track of which graphs a proof covers is part of keeping the statement contract of Habit~2 (Section~\ref{sec:magnitude}): the hypothesis ``acyclic'' belongs to the contract of Proposition~\ref{ch07:prop:acyclic-seymour}, and removing it changes the statement into the conjecture itself.
